% Master Thesis Report
% Created by: Claes Arvidson and Emelie Karlsson 
% Linköping University Spring 2016
% First version: 2016-01-27

%            the super-template for a thesis at LiTH v 1.002
%            ===============================================
%
% This is a template for writing a thesis at Linköpings Universitet
%
%   1 If you use the template to produce a final thesís (or some
%     other document) consider the template to be in the public domain.
%     This means that you can use and exploit it in any way. Otherwise
%     read 2 below.
%
%   2 This template is licenced under a Creative Commons Attribution, 
%     ShareAlike 2.5 license meaning that: 
%
%     You are free to:
%       1. Share: to copy, distribute, display, and perform the work.
%       2. Remix: to make derivative works.
%
%     Under the following conditions:
%       1. Attribution. You must attribute the work in the manner
%          specified by the author or licensor (see 11 below).
%       2. Share Alike. If you alter, transform, or build upon this
%          work, you may distribute the resulting work only under
%          a license identical to this one.  t
%       3. For any reuse or distribution, you must make clear to
%          others the license terms of this work.
%       4. Any of these conditions can be waived if you get permission
%          from the copyright holder (see 1 above).
%
%     (See also http://creativecommons.org/licenses/by-sa/2.5/ )
%
%   3 Example: if you use this template to produce another template
%     you must license it under the above license. And leave the 
%     changelog intact (you may or may not add your name to it).
%    
%     Example: if you use this template to produce a final thesis
%     you do not need to license your final thesis in any particular
%     way - but if you find significant improvements to this template
%     you may or may not share them.
%
%   4 if you find it useful - use it
%
%   5 if you do not find it useful - only use it for inspiration :)
%
%   6 if you find obvious errors: mail me or change them your self.
%
%   7 The text here are comments, when compiling this document
%     lines starting with % will be ignored. 
%
%   8 You can read the comments then erase them and save the file as
%     my_favourite_filename.tex to save the comments for a rainy day.
%
%   9 I have tried to make the strange pages automatic so that you do 
%     not have to edit them by hand.
%
%  10 I have tried to explain tables, figures etc "on the fly" so that
%     it will be easy for you to consult this template 
%
%  11 This template was compiled by me Per Erik Strandberg,
%     www.pererikstrandberg.se while I wrote my D-uppsats in applied
%     mathematics 2003/2004 at LiTH (accually it was finished after I
%     had finished it), since I felt the need of a simple template. 
%
%  12 If you make contributions to this template and redistribute it
%     feel free to add your name to the changelog below, but leave 
%     all previous authors (this is required by the license).
%
%
% Changelog (Authorname, email, date, contribution) 
% -----------------------------------------------------
%   * Per Erik Strandberg, per9000@gmail.com, 2004 apr 03: 
%     From different sources added 2 frontpages, the library page,
%     contents and the copyrightpage into one simple overviewable file.
%     Added comments and (semi-)automatic commands for boring stuff to
%     fill in manually (after having erased errorous automatic commands)
%
%   * Per Erik Strandberg, per9000@gmail.com, 2004 apr 30:
%     Some spelling problems were fixed.
%
%   * Per Erik Strandberg, per9000@gmail.com, 2007 jan 23: released to
%     the public domain for use only, redistribution of the template 
%     licensed under Creative Commons Attribution-ShareAlike 2.5.
%
%
% Stuff included are:
% -------------------
%   * Frontpage (fully automatic)
%   * another frontpage (fully automatic)
%   * library page (semiautomatic)
%   * contents (preface, chapters, bibliography, appendix)
%   * copyrightpage
%   * my walkthrough, and a checklist :)
%   * how to make an index (not that detailed)
%
% Files needed for this template:
% -------------------------------
%   * this file: template.tex
%   * some LaTeX (and TeX?)-packages provided by administrator or downloaded
%     to install them ask an administrator or use Miktex for windows.
%   * some pictures in eps-format 
%        - univlogo.eps   202 kb
%        - epresslogo.eps 190 kb
%
% Recommended guides:
% -------------------
%   * the not so short introduction to LaTeX available at:
%     http://www.ctan.org/tex-archive/info/lshort/english/lshort.pdf
%     or google for "the not so short introduction to LaTeX"
%   * your favourite local guide to LaTeX
%
% CHECKLIST - what MUST be done
% -----------------------------
%   * Fill in the newcommands starting on line ~ 80: author, abstract etc
%   * If you are not writing a thesis on D level for 20 credits 
%     you must change them to C and 10 or whatever on line ~ 200
%   * The library page is automatic except for
%        - department adress, line ~ 248
%        - language, and report category (well marked)
%        - and ISSN (0348-2960), line ~ 433 (just put nothing there if
%          you do not know what it is or if you are not wrinting a
%          thesis at MAI)
%
% KNOWN PROBLEMS AND SOLUTIONS:
% ------------------------------
%   P When in 12 pt mode the word licentiatavhandling is too large for
%     the library page, it looks nice in 10 pt mode however.  
%   S I just ignored this and hoped no one would mind.
% 
%   P Are the titlepages really centered in a good way?
%   S I don't know- you tell me.

% comments will be added in "realtime" like this:

% in documentclass tag you could write 10pt for a more printerfriendly format,
% or, if it is 10pt already, you could write 12pt for a more printer-
% -unfriendly format working better in S5-size (the bookletsize).
% But do not change a4paper to something else, the printing people
% will just rescale your document to the appropriate size.
% openright will force new chapter to start on the right
\documentclass[a4paper, 10pt, twoside, openright]{book}
\pagestyle{headings}

% some of the packages used are
%\usepackage[utf8]{inputenc}
\usepackage{latexsym}
\usepackage{eepic}
\usepackage{makeidx}
\usepackage[dvips]{graphicx}
\usepackage[utf8]{inputenc}  %för ä ö å ?
\usepackage[english]{babel}
%\usepackage{times}
\usepackage{amssymb}
\usepackage{fancybox} 
\usepackage{textcomp}

% the fancy header/footer
% consult http://research.cmis.csiro.au/gjw/tex/docs/fancyhdr.pdf
% or some other fancy header documenatation for more info
\usepackage{fancyhdr}
\pagestyle{fancy}
\lfoot{\nouppercase{\rightmark}}
\rfoot{\nouppercase{\leftmark}}
\fancyfoot{} 
\fancyhead[LO]{\nouppercase{\small{\rightmark}}}
\fancyhead[RE]{\nouppercase{\small{\leftmark}}}
\fancyhead[LE,RO]{\small{\thepage}}
\fancyhead[C]{}
\renewcommand{\footrulewidth}{0pt}
\renewcommand{\headrulewidth}{0.2pt}
\fancypagestyle{plain}{ \fancyhf{}  
  \fancyfoot[RE,LO]{\nouppercase{\small{\putegotext}}}
  \fancyfoot[LE,RO]{\small{\thepage}}
  \renewcommand{\headrulewidth}{0pt}  \renewcommand{\footrulewidth}{0.2pt} }

% ------------------------------------------------------------------------
% BEGINNING of abstract, author, examiner, etc
% here you must fill in your own text. Remember to respect the { and the }.
%
% the *-marked % in 
%                                   *
%      \newcommand{\putabstract}[0]{% 
%      some text}
%
% is made to avoid a blank space in the beginning of the abstract,
% author, date etc.
%
\newcommand{\putabstract}[0]{%
Here is where you can write your abstract. It may be very long, or it
may be very short, the reason you have an abstract is for people not
to be forced to read lots of crap.

But still, they will have to read your abstract. After all, the
abstract is what everyone reads\ldots}

\newcommand{\putkeywords}[0]{%
Keyword One, Chemostat, Another Key-Word, Key, Clé, Mot de cle,
Nyckelhål, XBOX, Dagens viktigaste nyckelord, and Keywords.}


\newcommand{\putthemonth}[0]{June }%please keep the space after the month

\newcommand{\putshortdate}[0]{2016}

% you don't have to change this
\newcommand{\putmydate}[0]{\putthemonth \putshortdate}

% this construction might be bad if you are multiple authors, 
% then comment the 4 lines below and
% uncomment the 2 lines below them 
%\newcommand{\putfirstname}[0]{Per Erik } % please keep the space
%
%\newcommand{\putlastname}[0]{Strandberg}
%
%\newcommand{\putauthor}[0]{\putfirstname \putlastname}
%
%\newcommand{\putegotext}[0]{\putlastname, \putshortdate.}

\newcommand{\putauthor}[0]{Claes Arvidson, Emelie Karlsson}

\newcommand{\putegotext}[0]{Arvidson, Karlsson, 2016.}


\newcommand{\puttitle}[0]{Scheduling of library staff with respect to demands and staff competence.}

% some kind of administrator will give you this number
\newcommand{\putregnumber}[0]{LiTH - MAT - EX - - 04 / 04 - - SE}

\newcommand{\putexaminer}[0]{E. Rönnberg}

\newcommand{\putsupervisor}[0]{T. Larsson}

% for example Applied Mathematics
\newcommand{\putdepartment}[0]{Very Applied Mathematics}

% send an email to ep.liu.se to make sure it is correct before you print
% this number is the one I had
\newcommand{\putmyurl}[0]{http://www.ep.liu.se/exjobb/mai/2004/tm/004/}

% Please do not change into University of Linköping or something, the
% international name of Liu is Linköpings Universitet
\newcommand{\putliu}[0]{Linköpings Universitet}

%
% END of abstract, author, examiner, etc
% ------------------------------------------------------------------------

\begin{document} 

\frontmatter

% some new commands I used, you won't need them I guess- just erase it
\newcommand{\ala}[0]{{\alpha}_{1}}
\newcommand{\alb}[0]{{\alpha}_{2}}
\newcommand{\nc}[0]{{N}_{C}}
\newcommand{\tCh}[0]{the Chemostat }
\newcommand{\Ch}[0]{Chemostat }
\newcommand{\ch}[0]{chemostat }
\newcommand{\MMk}[0]{Michaelis-Menten kinetics }
\newcommand{\KI}[0]{{K}_{I} }
\newcommand{\KN}[0]{{K}_{N} }
\newcommand{\mmu}[0]{{\mu}_{max} }
\newcommand{\KMAX}[0]{\mmu}
\newcommand{\AAA}[0]{\mathbf{A} }
\newcommand{\enc}[0]{utf8}
% stop erasing here

% don't mind these
\title{\puttitle}
\author{\putauthor}

% this command makes the first page, and the second page, handle with care
% You really don't have to change anything here I guess.
\renewcommand{\maketitle}{%
              \newpage%
              \pagestyle{empty}
              \null%
              \hfil\hspace*{-3mm}
              \begin{minipage}{150mm}
                \center 
                {\Large{Examensarbete}} \\\vspace*{4mm}
                {\vbox to 22mm{\vfil\Large\textbf{\puttitle}}}
                \vspace*{5mm}
                 {\Large{\putauthor}}       \\ \vspace*{4mm}
                  \putregnumber\\
                \end{minipage}\hfil
              \clearpage%\hspace*{0mm}\thispagestyle{empty}
              \clearpage%\thispagestyle{empty}
              \null%
              \hfil\hspace*{-3mm}
              \begin{minipage}{150mm}
                \center
                {\vbox to 48mm{\vfil\Large\textbf{\puttitle}}}
                \vspace*{5mm}
                  \putdepartment, \putliu \\ \vspace*{4mm}
                  \textbf{\putauthor}       \\ \vspace*{4mm}
                  \putregnumber
                  \vspace*{100mm}\\
                \flushleft
                
                Examensarbete:\hspace*{3pt}
                \begin{minipage}[t]{70mm}
% but you might want to change this 30
                  \textbf{30 hp}
                \end{minipage} \\ \vspace*{4mm}
                
                Level:\hspace*{3pt}
                \begin{minipage}[t]{70mm}
% and you might want to change this A 
                  \textbf{A} 
                \end{minipage} \\ \vspace*{4mm}
                
                Supervisor:\hspace*{3pt}
                \begin{minipage}[t]{120mm}
                  \textbf{\putsupervisor}, \\\putdepartment, \putliu
                \end{minipage} \\ \vspace*{4mm}
                
                Examiner:\hspace*{3pt}
                \begin{minipage}[t]{120mm}
                  \textbf{\putexaminer}, \\\putdepartment, \putliu
                \end{minipage} \\ \vspace*{4mm}
                
                Linköping:
                \begin{minipage}[t]{70mm}
                  \textbf{\putmydate}
                \end{minipage} \\ \vspace*{4mm}
              \end{minipage} \\ \hfill
}

% now the \maketitle command is changed and we use the new one
\maketitle

% I'm not sure this phantom is needed, but I leave to make sure
% a phantom leaves empty space corresponding to the size of what is in it.
\phantom{crap}
\thispagestyle{empty}
\pagestyle{empty}

% theese are old stuff made by the old library page, not needed I guess
%\setlength{\parindent}{0pt}
%\setlength{\textwidth}{210mm}
%\setlength{\textheight}{297mm}
%\setlength{\topmargin}{-1.0in}
%\setlength{\oddsidemargin}{-1.2in}
\setlength{\unitlength}{1mm}
%\begin{picture}(210,297)(40,-40)



\setlength{\unitlength}{1pt}

% now turn on fancy headers
\pagestyle{fancy}

% ..................chapter ...............................................
\chapter*{Abstract\index{Abstract}}
% your abstract
\putabstract

% enters keywords
\begin{description}
\item[Keywords:]{%
\putkeywords
}
\item[URL for electronic version: ]{\hfill%\quad\\%
%\phantom{\qquad} http://urn.kb.se/resolve?urn=urn:nbn:se:liu:diva-77777\\
%where 77777 should be replacd by an appropriate number.
\begin{center}
http://urn.kb.se/resolve?urn=urn:nbn:se:liu:diva-77777
\end{center}
%where 77777 should be replaced by an appropriate number.
}
\end{description}


% here or on the next page you put an abstract in swedish
%\section*{Abstract in Swedish: Sammanfattning}
%Det finns många typer av bioreaktorer som tillämpas\ldots

\chapter*{Acknowledgements\index{Acknowledgements}}

I would like to thank my supervisor, I would like to thank my
supervisor, I would like to thank my supervisor, I would like to thank my supervisor\ldots

I also have to thank, I would like to thank my supervisor, I would
like to thank my supervisor, I would like to thank my supervisor, I
would like to thank my supervisor\ldots

My opponent NN also deserves my thanks, I would like to thank my
supervisor, I would like to thank my supervisor, I would like to thank my supervisor\ldots





% optional I guess
\chapter*{Nomenclature\index{Nomenclature}}

Most of the reoccurring abbreviations and symbols are described here.

%----------------NEW_SECTION--------------------------------------------------

\section*{Symbols\index{Symbols}}

\begin{tabular}{ll}
$Y_0$      & The amount of the variable $Y$ inserted into a system.\\
$\hat Y$& The unit-dimension of the variable $Y$, for example $\hat t=1s$ .\\
$\bar Y_i$ & A steady state (number $i$) value of Y.\\
\phantom{a}& \phantom{b} \\
$K_i$ & Constants used in kinetic expressions, for example $K_I$.\\
\phantom{a}& \phantom{b} \\
$\AAA$     & The system matrix. \\
\end{tabular}

\section*{Abbreviations\index{Abbreviations}}

\begin{tabular}{ll}
CPI      & Competitive Product Inhibition (or Inhibited)\index{CPI}\\
CSI      & Competitive Substrate Inhibition (or Inhibited)\index{CSI}\\
CSTR     & Continuous Stirred Tank (bio)Reactor\index{CSTR}\\
MMI      & Michaelis-Menten Inhibition (or Inhibited)\index{MMI}
\end{tabular}


\tableofcontents 

\mainmatter


% ..................chapter ...............................................
\chapter{Introduction}

\iffalse

% I wanted sans serif in the ingress of every chapter, some people say
% this is possible with the command \ingress or something, I don't know...
\textsf{%
This text is written as a master of science final
thesis\footnote{No it is not- it is a LaTeX-template for people
  writing them.}
at Linköpings universitet
by Per Erik Strandberg
% sometimes you need ~ to force space
with \putsupervisor~as supervisor and \putexaminer~as examiner,
with Emacs in \LaTeX{}
in 2003/2004.}

\textsf{%
This first chapter will define what is a Continuous stirred tank bioreactor
and a chemostat. We give some background and formulate some main
questions we wish to answer, and describe what topics are covered.  
}


\section{Background} 

Modeling of biological growth in reactors\ldots
by dynamical systems started in the 1950's,
when Monod, and Novick and Szilard developed the concept of
\emph{Continuous Stirred Tank (bio)Reactors}(CSTRs\index{CSTR}) that
differed from traditional batchreactors.

% this is one example of how to make a picture
\begin{figure}[h]
\begin{center}
\begin{tabular}{|p{75mm}|}
\hline \\
\begin{picture}(130,090)
\put(080,020){\line(1,0){050}}
\put(080,020){\line(0,1){030}}
\put(130,020){\line(0,1){030}}
\multiput(080,045)(002,000){25}{\line(1,0){001}}
\put(090,085){\vector(0,-1){35}}
\put(123,025){\vector(1,0){25}}
\put(082,025){$X$, $S$, $\mu$}
\put(140,012){$F$, $X$, $S$}
\put(092,065){$S_0$, $F$}
\end{picture}\\
\hline
\end{tabular}
\end{center}
\caption{The ideal CSTR: the chemostat\ldots}
\end{figure}

A chemostat is typically made of a reactor\ldots 

\section{Another section}

This report provides a summary of already
existing knowledge of bioreactors and a detailed analysis\ldots


\section{Topics covered}
There are five chapters (and this introduction)
and two appendixes. Main topics dealt with are: 

% this is my favourite command: description...
\begin{description} 
\item[Chapter 2:]{We explain the chemostat equations with the standard
    Michaelis-Menten kinetics by starting with two simple models of
    biological growth. We also present some tools useful for studying
    dynamical systems in two dimensions.}   
\item[Appendix A:]{An appendix dealing with tools used to simplify the
    analysis of linearized stability in steady states.}
\end{description}

% -------------------------------------------------------------------------
\fi
% ----------------------------------From here ------------------------------

\inputencoding{\enc}
% Introduction to Master Thesis

\section{Background}


At a library, staff absence can cause problems and result in a shortage of staff with a special competence. If a worker is unavailable a certain day due to a meeting or illness it would require a stand-in to fill the vacancy. Therefore, it is of great interest to create schedules with as many skilled stand-ins as possible to overcome such disturbances, when they occur. Furthermore, the library personnel have certain demands and preferences as to how a satisfactory working schedule should be. For instance, it is neither preferable to work more than one evening per week nor work more weekends than required. 

The problem addressed in this thesis work concerns the library staff at the Central Library of Norrköping. The library has more than 40 employees and the culturally important building from the 1950's is known to all inhabitants of Norrköping. The library is open weekdays between 08:00 and 20:00 and during weekends between 11:00 and 16:00. The generous opening times also creates a challenge for the library as it requires a large pool of well coordinated personnel to keep it running. In addition, the library also provides its services to one other smaller library. This creates a challenge for the library schedulers.

\section{Problem description}

In this section, the most important features of the worker scheduling problem studied in this thesis work will be explained.  

\subsection{Description of the daily tasks at the library} \label{section:library_tasks}
The most important activity at a library is the activity directed towards the public. This includes book loan services as well as providing customers with helpful information about the resources at the disposal of the library. These activities are referred to as "outer tasks". In addition, the book collection must be maintained, the returned books must be sorted and put back to the shelves, the web page must be updated and so on. Such work is often referred to as "inner work" and this non-public work is an equally important part of the everyday activities at the library. 

Three main outer task types can be identified at the library of Norrköping, working in the service counter (sv. expiditionsdisken), working in the information counter (sv. informationsdiken)  and assembling books which are to be sent to other libraries according to the "fetch list" (sv. plocklista). The fetch list is a task type for which the worker is scheduled during a whole day, while the other task types have two to five hour shifts. The outer tasks can be performed by either librarians or both, as described in Table \ref{tab:Outer_Tasks}.

\begin{table}[h]
\centering
\caption{Outer tasks can be performed exclusively by librarians or by both librarians and assistants.}
\label{tab:Outer_Tasks}
\begin{tabularx}{\textwidth}{|l|l|X|}
\hline
%-------------------------------------------------------------------
\textbf{Task} & \textbf{Description} & \textbf{Qualification}\\ \hline 
%------------------------------------------------------------------- 
\specialcell[t]{Service counter \\ (Exp)}  & \specialcell[t]{Administring loans, library cards\\ and the loaning machine} & \specialcell[t]{Librarian or \\  assistant} 
%\begin{tabular}[x]{@{}c@{}}\\\end{tabular}  
\\ \hline
%-------------------------------------------------------------------
\specialcell[t]{Information counter \\ (Info)} & \specialcell[t]{Handling questions \\about the library's resources.} & Librarian
\\ \hline 
%-------------------------------------------------------------------
\specialcell[t]{Fetch list \\(PL)} & \specialcell[t]{Fetching books that are to be \\sent to other libraries.} & \specialcell[t]{Librarian  or \\  assistant}
\\ \hline 
%-------------------------------------------------------------------
\specialcell[t]{Hageby \\(HB)} & \specialcell[t]{Handling librarian tasks at the filial \\ Hageby during weekends.} & Librarian
\\ \hline 
%-------------------------------------------------------------------
\specialcell[t]{Library on Wheels \\(BokB)} & \specialcell[t]{Driving the Library on Wheels \\ to different areas of town.} & Librarian
\\ \hline 
%-------------------------------------------------------------------
\end{tabularx}
\end{table} 

Since the number of visitors at the library varies throughout the day and also during different days of the week, so does the demand for personnel for the three tasks. The demand of personnel for the different task types is illustrated in Table \ref{tab:Outer_Task_Demand}, according to figures given by the library. 



\begin{table}[h]
\centering
\caption{Demand of staff for the three daily tasks.}
\label{tab:Outer_Task_Demand}
\begin{tabularx}{\textwidth}{|X|l|l|l|l|X|}
\hline
%-------------------------------------------------------------------
\textbf{Day} & \textbf{Time} & \textbf{Exp demand} & \textbf{Info demand} & \textbf{PL demand\footnote{The Fetch list is a whole day task covering three consecutive shifts which must be performed by the same worker.}}
\\ \hline 
%------------------------------------------------------------------- 
Mon-Fri & 08:00-10:00 & 2 & 2 & 1
\\ \hline 
%------------------------------------------------------------------- 
Mon-Fri & 10:00-13:00 & 3 & 3 & 1
\\ \hline 
%-------------------------------------------------------------------
Mon-Fri & 13:00-16:00  & 3 & 3 & 1
\\ \hline 
%-------------------------------------------------------------------
Mon-Fri & 16:00-20:00 & 3 & 3 & -
\\ \hline 
%-------------------------------------------------------------------
Sat & 11:00-16:00  & 3 & 3 & -
\\ \hline 
%------------------------------------------------------------------- 
Sun & 11:00-16:00  & 3 & 3 & -
\\ \hline 
%------------------------------------------------------------------- 
\end{tabularx}

\end{table} 

As in the case with most libraries, the Central Library of Norrköping also has responsibilities that fall outside of its normal daily activities. One such responsibility is the running of a smaller library filial in Hageby, situated in a suburban area of Norrköping, during weekends. It is decided that only librarians are qualified for this task, as it implies some tasks that need certain knowledge.

Table \ref{tab:Hageby_Demand} shows the demand of staff at Hageby. Worth noting is that it is the same person working at Hageby a weekend due to a couple of reasons. Since a worker is supposed to work both Saturday and Sunday when due for a weekend, it is more desirable to let the worker focus on the same task both days. Furthermore, since Hageby is located in a suburban, some travel distance is required to reach it. This is compensated by letting the worker at Hageby be free from work Friday evening, which otherwise is included in the weekend work. 

\begin{table}[H]
\centering
\caption{Demand of staff at Hageby}
\label{tab:Hageby_Demand}
\begin{tabularx}{0.5\textwidth}{|l|l|X|}
\hline
%-------------------------------------------------------------------
\textbf{Day} & \textbf{Time} & \textbf{HB demand}
\\ \hline 
%%------------------------------------------------------------------- 
Sat & 11:00-16:00  & 1 
\\ \hline 
%%-------------------------------------------------------------------
Sun & 11:00-16:00  & 1 
\\ \hline 
\end{tabularx}
\end{table} 

Similarly, only a handful librarians are qualified for the task known as the "Library on Wheels" (sv. Bokbussen), which is a library bus that provides citizens in remoter areas of the city with books and other library services. The Library on Wheels only operates a few times a week and the schedule differs between odd and even weeks.

\begin{table}[!h]
\centering
\caption{Demand of staff Library on Wheels}
\label{tab:LOW_Demand}
\begin{tabularx}{0.80\textwidth}{|l|X|X|X|X|X|}
\hline
%%-------------------------------------------------------------------
 \textbf{Odd Week} & \textbf{Mon} & \textbf{Tue} & \textbf{Wed} & \textbf{Thu} & \textbf{Fri} 
 \\ \hline 
%%------------------------------------------------------------------- 
\rowcolor{Gray} 
08:00-10:00 & 1 & 0 & 1 & 1 & 1 
\\ \hline 
%%------------------------------------------------------------------- 
\rowcolor{Gray} 
16:00-20:00 & 1 & 0 & 1 & 1 & 0 
\\ \hline 
%%-------------------------------------------------------------------
 \textbf{Even Week} & \textbf{Mon} & \textbf{Tue} & \textbf{Wed} & \textbf{Thu} & \textbf{Fri} 
 \\ \hline 
%%------------------------------------------------------------------- 
\rowcolor{Gray} 
08:00-10:00 & 1 & 0 & 1 & 1 & 0 
\\ \hline 
%%------------------------------------------------------------------- 
\rowcolor{Gray} 
16:00-20:00 & 0 & 0 & 1 & 1 & 0 
\\ \hline 
\end{tabularx}
\end{table}

Apart from the outer work described above, there is also inner work at the library which sometimes need to be scheduled. One such type of inner work is meetings, as it concerns a large number of staff. Both library meetings, scheduled for the whole work force, and department meetings, scheduled only a subset of workers, exist. The library meetings are fixed in time and take place every fifth weeks on Monday mornings 8:15-10:00 while the department meetings are more flexible and take place every fifth week. For this, all department members must be available. Only three departments have scheduled meetings: the child department, the adult department and the media department.

\subsection{Personnel attributes}

At a library, the main group of workers are librarians and assistants (also janitors, cleaners, security guards etc.). The librarians make sure the library's resource collection is up to date, that the visitors find what they're looking for and perform a larger number of inner tasks. Assistants can perform most of theses tasks, while others require librarian competence.

The Central Library of Norrköping currently has 39 workers, 23 of which are librarians and 16 of which are assistants. All staff have different availability for performing tasks, depending on their working hours and the amount of inner work they are in charge of. In the standard case, each worker is assigned one evening per week and once per five weeks he/she is assigned to work during the weekend. The following week after a weekend is compensated by two free days, placed according to the wishes of the worker.

Let us consider a sample worker who is qualified for the same tasks as a librarian, works full time and is also assigned to work on Wednesday evenings. The worker is also assigned to work weekend on the fourth week and has chosen to take out its days off on Thursdays and Fridays on week five. The workers availability is thus illustrated in table \ref{tab:Bob_avail}. The schedule repeats itself after five weeks and illustrates only the availability for tasks but says nothing about whether the worker has been assigned any tasks or not.

All workers have a five week schedule in the same manner as the sample worker. However, in order to meet the weekend demands illustrated by Tables \ref{tab:Outer_Task_Demand}, \ref{tab:LOW_Demand} and \ref{tab:Hageby_Demand}, it is evident that not all workers can be assigned for weekend work at week 4. Instead, a worker's schedule should be seen as a \textit{relative schedule}, which can be shifted by whole weeks. We refer to this as the \textit{week rotation}.

The relative schedule is relative to the \textit{overall schedule} for the library. If the sample worker for example has week rotation 1, this means the weekend week in his relative schedule would be shifted to week 1 in the overall schedule. In Table \ref{tab:Bob_avail}, the weekend rotation is 4.



\begin{table}[!h]
\centering
\caption{Availability schedule for a sample worker. Yellow signifies that the worker is available. In parenthesis, the weekend shift. The week rotation is 4}
\label{tab:Bob_avail}
\begin{tabularx}{\textwidth}{|X|l|l|l|l|l|l|l|X|}
\hline
%-------------------------------------------------------------------
\textbf{Week 1}& \colcell \textbf{Mon} & \colcell \textbf{Tue} & \colcell \textbf{Wed} & \colcell \textbf{Thu} & \colcell \textbf{Fri} & \colcell \textbf{Sat} & \colcell \textbf{Sun}
\\ \hline 
%%------------------------------------------------------------------- 
%\rowcolor{Gray} 
\colcell 08:00-10:00 (11:00-16:00) & \colcelltwo & \colcelltwo & \colcelltwo & \colcelltwo & \colcelltwo & & 
\\ \hline 
%%-------------------------------------------------------------------
%\rowcolor{Gray} 
\colcell 10:00-13:00 & \colcelltwo & \colcelltwo & \colcelltwo & \colcelltwo & \colcelltwo &   & 
\\ \hline 
%%-------------------------------------------------------------------
%\rowcolor{Gray} 
\colcell 13:00-16:00 & \colcelltwo & \colcelltwo & \colcelltwo & \colcelltwo & \colcelltwo & &
\\ \hline 
%%-------------------------------------------------------------------
%\rowcolor{Gray} 
\colcell 16:00-20:00 & & & \colcelltwo & & & &
\\ \hline 
%%-------------------------------------------------------------------
\end{tabularx}
\begin{tabularx}{\textwidth}{|X|l|l|l|l|l|l|l|X|}
\hline
%-------------------------------------------------------------------
\textbf{Week 2}& \colcell \textbf{Mon} & \colcell \textbf{Tue} & \colcell \textbf{Wed} & \colcell \textbf{Thu} & \colcell \textbf{Fri} & \colcell \textbf{Sat} & \colcell \textbf{Sun}
\\ \hline 
%%------------------------------------------------------------------- 
%\rowcolor{Gray} 
\colcell 08:00-10:00 (11:00-16:00) & \colcelltwo & \colcelltwo & \colcelltwo & \colcelltwo & \colcelltwo & & 
\\ \hline 
%%-------------------------------------------------------------------
%\rowcolor{Gray} 
\colcell 10:00-13:00 & \colcelltwo & \colcelltwo & \colcelltwo & \colcelltwo & \colcelltwo &   & 
\\ \hline 
%%-------------------------------------------------------------------
%\rowcolor{Gray} 
\colcell 13:00-16:00 & \colcelltwo & \colcelltwo & \colcelltwo & \colcelltwo & \colcelltwo & &
\\ \hline 
%%-------------------------------------------------------------------
%\rowcolor{Gray} 
\colcell 16:00-20:00 & & & \colcelltwo & & & &
\\ \hline 
%%-------------------------------------------------------------------
\end{tabularx}
\begin{tabularx}{\textwidth}{|X|l|l|l|l|l|l|l|X|}
\hline
%-------------------------------------------------------------------
\textbf{Week 3}& \colcell \textbf{Mon} & \colcell \textbf{Tue} & \colcell \textbf{Wed} & \colcell \textbf{Thu} & \colcell \textbf{Fri} & \colcell \textbf{Sat} & \colcell \textbf{Sun}
\\ \hline 
%%------------------------------------------------------------------- 
%\rowcolor{Gray} 
\colcell 08:00-10:00 (11:00-16:00) & \colcelltwo & \colcelltwo & \colcelltwo & \colcelltwo & \colcelltwo & & 
\\ \hline 
%%-------------------------------------------------------------------
%\rowcolor{Gray} 
\colcell 10:00-13:00 & \colcelltwo & \colcelltwo & \colcelltwo & \colcelltwo & \colcelltwo &   & 
\\ \hline 
%%-------------------------------------------------------------------
%\rowcolor{Gray} 
\colcell 13:00-16:00 & \colcelltwo & \colcelltwo & \colcelltwo & \colcelltwo & \colcelltwo & &
\\ \hline 
%%-------------------------------------------------------------------
%\rowcolor{Gray} 
\colcell 16:00-20:00 & & & \colcelltwo & & & &
\\ \hline 
%%-------------------------------------------------------------------
\end{tabularx}
\begin{tabularx}{\textwidth}{|X|l|l|l|l|l|l|l|X|}
\hline
%-------------------------------------------------------------------
\textbf{Week 4}& \colcell \textbf{Mon} & \colcell \textbf{Tue} & \colcell \textbf{Wed} & \colcell \textbf{Thu} & \colcell \textbf{Fri} & \colcell \textbf{Sat} & \colcell \textbf{Sun}
\\ \hline 
%%------------------------------------------------------------------- 
%\rowcolor{Gray} 
\colcell 08:00-10:00 (11:00-16:00) & \colcelltwo & \colcelltwo & \colcelltwo & \colcelltwo & \colcelltwo & \colcelltwo & \colcelltwo
\\ \hline 
%%-------------------------------------------------------------------
%\rowcolor{Gray} 
\colcell 10:00-13:00 & \colcelltwo & \colcelltwo & \colcelltwo & \colcelltwo & \colcelltwo &   & 
\\ \hline 
%%-------------------------------------------------------------------
%\rowcolor{Gray} 
\colcell 13:00-16:00 & \colcelltwo & \colcelltwo & \colcelltwo & \colcelltwo & \colcelltwo & &
\\ \hline 
%%-------------------------------------------------------------------
%\rowcolor{Gray} 
\colcell 16:00-20:00 & & & \colcelltwo & & \colcelltwo & &
\\ \hline 
%%-------------------------------------------------------------------
\end{tabularx}

\begin{tabularx}{\textwidth}{|X|l|l|l|l|l|l|l|X|}
\hline
%-------------------------------------------------------------------
\textbf{Week 5}& \colcell \textbf{Mon} & \colcell \textbf{Tue} & \colcell \textbf{Wed} & \colcell \textbf{Thu} & \colcell \textbf{Fri} & \colcell \textbf{Sat} & \colcell \textbf{Sun}
\\ \hline 
%%------------------------------------------------------------------- 
%\rowcolor{Gray} 
\colcell 08:00-10:00 (11:00-16:00) & \colcelltwo & \colcelltwo & \colcelltwo & & & & 
\\ \hline 
%%-------------------------------------------------------------------
%\rowcolor{Gray} 
\colcell 10:00-13:00 & \colcelltwo & \colcelltwo & \colcelltwo & & & & 
\\ \hline 
%%-------------------------------------------------------------------
%\rowcolor{Gray} 
\colcell 13:00-16:00 & \colcelltwo & \colcelltwo & \colcelltwo & & & &
\\ \hline 
%%-------------------------------------------------------------------
%\rowcolor{Gray} 
\colcell 16:00-20:00 & & & \colcelltwo & & & &
\\ \hline 
%%-------------------------------------------------------------------
\end{tabularx}
\end{table} 

Considering again the sample worker's schedule, it may look as if the worker can be scheduled for tasks at any yellow shift. However, there are regulations set by the library controlling how the outer tasks are to be distributed among the workers. The purpose of the regulations is to guarantee that the worker schedules are not too unbalanced or uncomfortable. 

There four basic demands of the resulting schedule. Firstly, workers are only allowed to take at a maximum one task per day. Secondly, workers only work at most weekend per five weeks. During this weekend it is customary that you work Friday evening shift, Saturday and Sunday shift consecutively, unless you are scheduled for Hageby, which is far away and is thus compensated by not entailing Friday evening work. Thirdly, weekend work is to be compensated by days off. There are a few possible variations of the weekend compensation, but usually the worker takes one and a half or two days off the week following upon weekend work. Lastly, a worker is only allowed to work one evening per week, excluding the Friday which belongs to the weekend week.
%
%Regulations say that:
%
%\begin{enumerate}
%\item A worker can at a maximum perform one task per day.
%\item A worker works only one weekend in five weeks.
%\item A worker must work Friday evening shift, Saturday shift and Sunday shift consecutively. For HB-workers, no Friday evening shift.
%\item A worker is only allowed to work one evening per week, excluding the Friday evening of the weekend week.
%\end{enumerate}

An example of a feasible schedule for the sample worker is provided in Table \ref{tab:Lib_feas_sched}. It should be noted that this schedule is created for a general worker and does not necessarily apply to all workers. Among workers, special availability restrictions due to odd-and even weeks exist and only a subset of the librarians are available for the Library on Wheels and Hageby. Also, some workers never work evening or weekend shifts.  

\begin{table}[!h]
\centering
\caption{Example of a feasible first week for the sample worker.}
\label{tab:Lib_feas_sched}
\begin{tabularx}{\textwidth}{|X|l|l|l|l|l|l|l|X|}
\hline
%-------------------------------------------------------------------
\textbf{Week 1} & \colcell \textbf{Mon} & \colcell \textbf{Tue} & \colcell \textbf{Wed} & \colcell \textbf{Thu} & \colcell \textbf{Fri} & \colcell \textbf{Sat} & \colcell \textbf{Sun}
\\ \hline 
%%------------------------------------------------------------------- 
%\rowcolor{Gray} 
\small \colcell 08:00-10:00 (11:00-16:00)& \colcelltwo & \small \colcellthree Exp & \colcelltwo & \small \colcellthree PL & \colcelltwo & & 
\\ \hline 
%%-------------------------------------------------------------------
%\rowcolor{Gray} 
\small \colcell 10:00-13:00 & \small \colcellthree Exp & \colcelltwo & \colcelltwo & \small \colcellthree PL & \colcelltwo & & 
\\ \hline 
%%-------------------------------------------------------------------
%\rowcolor{Gray} 
\small \colcell 13:00-16:00 & \colcelltwo & \colcelltwo & \colcelltwo & \small \colcellthree PL & \small \colcellthree Info & &
\\ \hline 
%%-------------------------------------------------------------------
%\rowcolor{Gray} 
\small \colcell 16:00-20:00 & & & \small \colcellthree Info& & & &
\\ \hline 
%%-------------------------------------------------------------------
\end{tabularx}
\end{table} 


\subsection{Scheduling objectives: stand-in maximization and schedule variation}

What is investigated in this thesis is not primarily how to distribute the tasks according to the rules given above, but rather how to do this is a way so that the number of stand-in personnel is maximized. The emergency back-up system of stand-ins is crucial for the library in order to be able to keep the library desks open and fully staffed. Thus, the highest priority is to maximize the number of stand-in personnel.

A stand-in is defined as a worker who is available for outer tasks during the first three shifts, as well as not scheduled for any shift that day. Inner work sometimes prevents personnel from being scheduled as a stand-in for the outer tasks, and therefore also information about such work must be known to the scheduler. Both librarians and assistants can be scheduled as a stand-in, but only librarians are qualified to take emergency shifts at the Info desk. Since the regular library activity is the most crucial activity, there is no need for stand-ins during evening and weekends. Similarly, there are no assigned stand-ins for the Library on Wheels or for Hageby.

Apart from maximizing the number of stand-ins for each day at the library, another important measurement of a good schedule is the level of variation. It is desirable for the sake of the personnel to create schedules in which the weeks resemble eachother or are repeated according to some pattern. On the other hand, since some shifts are more attractive than other and in order to avoid too much repetitiveness it is desirable that the days during a work week are not too similar. The example schedule in Table \ref{tab:Lib_feas_sched} is an example of a sufficiently varied weekly schedule. 

\section{Method}



\section{Topics Covered}













\iffalse



\begin{table}[h]
\centering
\caption{Overall schedule for week 1.}
\label{tab:General_schedule}
\begin{tabularx}{\textwidth}{|X|l|l|l|l|X|}
\hline
%-------------------------------------------------------------------
\textbf{Schedule week 1} & & & & &  
\\ \hline 
%%------------------------------------------------------------------- 
\textbf{Monday}& \colcell \textbf{Exp} & \colcell \textbf{Info} & \colcell \textbf{PL} & \colcell \textbf{LoW - odd} & \colcell \textbf{LoW- even} 
\\ \hline 
%%------------------------------------------------------------------- 

 08:00 - 10:00 & \small W1, W2 & \small W1,W2 & \small W1 & \small W1 & \small W1
\\ \hline 
%%------------------------------------------------------------------- 
 10:00 - 13:00 & \small W1, W2, W3 & \small W1, W2, W3 & \small W1 & - & -
\\ \hline 
%%------------------------------------------------------------------- 
 13:00 - 16:00 & \small W1, W2, W3 & \small W1, W2, W3 & \small W1 & - & -
\\ \hline 
%%------------------------------------------------------------------- 
 16:00 - 20:00 & \small W1, W2, W3 & \small W1, W2, W3 & \small W1 & - & -
\\ \hline 
%%------------------------------------------------------------------- 
\textbf{Tuesday}& \colcell \textbf{Exp} & \colcell \textbf{Info} & \colcell \textbf{PL} & \colcell \textbf{LoW - odd } & \colcell \textbf{LoW- even} 
\\ \hline 
%%------------------------------------------------------------------- 
\textbf{Wednesday}& \colcell \textbf{Exp} & \colcell \textbf{Info} & \colcell \textbf{PL} & \colcell \textbf{LoW - odd } & \colcell \textbf{LoW- even} 
\\ \hline 
%%-------------------------------------------------------------------
\textbf{Thursday}& \colcell \textbf{Exp} & \colcell \textbf{Info} & \colcell \textbf{PL} & \colcell \textbf{LoW - odd } & \colcell \textbf{LoW- even} 
\\ \hline 
%%-------------------------------------------------------------------
\textbf{Friday}& \colcell \textbf{Exp} & \colcell \textbf{Info} & \colcell \textbf{PL} & \colcell \textbf{LoW - odd } & \colcell \textbf{LoW- even} 
\\ \hline 
%%-------------------------------------------------------------------
\end{tabularx}
\end{table}

\fi

\inputencoding{\enc}

% ..................chapter ...............................................

\chapter{The ideal CSTR: the chemostat\index{Chemostat}}

\textsf{%
In this chapter we study exponential growth, the logistic\ldots .}

\section{Some simple models of biological growth}

\subsection{Exponential growth\index{Exponential growth}}

% the command mbox forces things to be placed together
If \mbox{$\mu =$ constant $>0$,} we get $X(t)={X}_0
{e}^{\mu t}$.

\subsection{The logistic equation\index{Logistic equation}}

Let us assume that $\frac{dX}{dt}=\mu \cdot X,$ with
$\mu=\mu(S)=k\cdot S$ \ldots 

% one way of displaying maths: 
\begin{displaymath}
\left\{ \begin{array}{ll}
\frac{dX}{dt}=kSX                            & (a) \\
\frac{dS}{dt}=- \alpha kSX \rule{0pt}{3ex}   & (b) 
\end{array}\right.
\end{displaymath}

% another way to display math with a number:
\begin{equation} \label{logistic}
\frac{dX}{dt}= r(1- \frac{X}{B})X  
\end{equation}

An explicit solution to (\ref{logistic}) is: 
$X(t)=\frac{{X}_{0}B}{{X}_{0}+(B-{X}_{0}){e}^{-rt}},$ if $0<{X}_{0}<B$. 
It can be found by separating variables in equation (\ref{logistic})


\section{The chemostat\index{Chemostat|textbf}} 

A chemostat is made of two main parts; a nutrient reservoir, and a
growth-chamber, reactor, in which the bacteria reproduces. 

% here is how to display text upright in formula \textrm{your text}
\begin{equation} \label{chemo_basic}
\left\{ \begin{array}{l}
\frac{dX}{dt}= \mu(S)X \overbrace{-X\frac{F}{V}}^{\textrm{new}}  \\
\frac{dS}{dt}=- \alpha \mu(S) X
\underbrace{-S\frac{F}{V}+ {S}_{0}\frac{F}{V}}_{\textrm{new}} 
\end{array}\right.
\end{equation}

% yet another way of displaying maths and a way of using a matrix
\begin{displaymath}
\mathbf{A} =
\left( \begin{array}{cc}
0 &  \sigma {\alpha}_{1}  \\
 - \frac{1}{{\alpha}_{1}} & -\sigma - 1 \\
\end{array}\right)
\end{displaymath}

% this subsubsection will not be numbered because we use the *
\subsubsection*{The invariant line: conclusions}




% the not so good structure of how a formula might look:
%\begin{displaymath}
%\mathbf{A}=
%\left( \begin{array}{cc}
%\ala \frac{S}{1+S} \frac{1}{(1+\frac{X}{\XC})^2} -1 
%& \ala \frac{1}{(1+S)^2} \frac{X}{1+\frac{X}{\XC}}   \\
%- \frac{S}{1+S} \frac{1}{(1+\frac{X}{\XC})^2} \phantom{+11} 
%& - \frac{1}{(1+S)^2} \frac{X}{1+\frac{X}{\XC}}-1  \rule{0pt}{3ex}   \\
%\end{array}\right).
%\end{displaymath}


% a big, bad table
\begin{tabular}{|c|c|c|} \hline %\hline
%-------------------------------------------------------------------
$\phantom{\Bigg(}$Model$\phantom{\Bigg(}$ & Monods Chemostat & CSI-CSTR \\ 
\hline %\hline 
%-------------------------------------------------------------------
$\mu\phantom{\Bigg(}$
& $\frac{ S} {1+S}$
& $\frac{ S} {1+S+\frac{{S}^{2}}{\KI}}$ \\ %\hline
%-------------------------------------------------------------------
$\frac{dX}{dt} \phantom{\Bigg(}$ 
& $\phantom{-}\ala\frac{ S}{1+S}X-X$        
& $\phantom{-}\ala\frac{ S}{1+S+\frac{{S}^{2}}{\KI}}X-X $  \\
$\frac{dS}{dt} \phantom{\Bigg)}$ 
& $-\frac{S} {1+S}X-S+\alb $        
& $-\frac{S} {1+S+\frac{{S}^{2}}{\KI}}X-S+\alb$        \\ %\hline
%------------------------------------------------------------------
XNC%
& $S=\frac{1}{\ala - 1} \phantom{\Bigg(}$
& $S=\frac{\KI(\ala-1)}{2}\pm\sqrt{{(\frac{\KI(\ala-1)}{2})}^{2}-\KI}$ \\
SNC% 
& $ X= \frac{ (\alb-S)(1+S) }  { S }$ 
& $ X=
\frac{({\alpha}_{2}-S)(1+S+\frac{{S}^{2}}{\KI})}{S}\phantom{\Bigg(} $\\ 
%-------------------------------------------------------------------
limit% 
& $-$
& $ K_I \rightarrow \infty\phantom{\Bigg(}$  \\ \hline 
\end{tabular}


% No matter how many empty lines you put here, it will be compact in
% your outfile. Hint: solve with some kind of phantom or so.


The other three models, the chemostat, the MMI-CSTR and the CPI-CSTR
are quite similar in comparison to the CSI-CSTR. 

% one way of using ``feel'' instead of "feel"
Monods chemostat does not ``feel'' this inhibition and does not care\ldots 


% this is a way of using the index subgroups
% \section{Chemostats in series\index{Chemostat!In series} 


% ..................chapter ...............................................

% you really should have some kind of summary and conclusion
% \chapter{Summary and conclusion}
% In this report several mathematical models of bacteria population
% growth in bioreactors were analyzed.


% .................. bibliography...................................
% the number 99 is some kind of maximum for LaTeX to know an upper limit
% this is one of the ways of dealing with references, when you refer
% to something here you will get: "...according to [1] we have..." or
% "...another observation made by [3,6,60] suggests that..."
%
Here is an example of how to cite books in your bibliography. This
text will be displayed at the end of chapter two. This is some kind of
bibiography, according to \cite{beta}, we have\ldots And according to
\cite{beta,Glad} we have something else.

\begin{thebibliography}{99}

\bibitem{beta} % this beta is what you 'call' when you refer to beta
Lennart Råde, Bertil Westergren, 
(2001), 
\emph{Mathematics Handbook for Science and Engineering},
Studentlitteratur, 
Lund.

\bibitem{Glad} % and glag is the handle of Glad.
Torkel Glad, Lennart Ljung,
(1989), 
\emph{Reglerteknik grundläggande teori},
Studentlitteratur, 
Lund.

\end{thebibliography}

% .................. appendix ...................................

\appendix

\chapter{The Linearized stability\index{Stability}}

\section{The Linearization\index{Linearization}}

$F(x)$, a one-variable function of $x$ can be Taylor-expanded
around a fix $X$. We get $F(X+x)=F(X) + F'(X) x + O({x}^{2})$.
For small perturbations of $x$ around $X$ we get the linearization:
$F(X+x) \approx F(X) + F'(X) x$, containing only the
constant and the linear terms.

For functions of two variables $F(X+x,S+s)$ and $G(X+x,S+s)$:
\begin{displaymath} 
\left\{ \begin{array}{l}
F(X+x,S+s)= F(X,S) + {F}_{X}'(X,S)x + {F}_{S}'(X,S)s + O({(x+s)}^{2})\\
G(X+x,S+s)= G(X,S) + {G}_{X}'(X,S)x + {G}_{S}'(X,S)s + O({(x+s)}^{2})
\rule{0pt}{3ex} 
\end{array}\right. 
\end{displaymath}

% I had some matlab code, insert it like this if you feel like it:
% scriptsize lowers fontsize
% verbatim forces LaTeX to write all symbols and in your way
\begin{scriptsize}
\begin{verbatim}
function chemostat_inhibited(alpha1, alpha2, xp0, sp0, xc)
%
%chemostat_inhibited Displays a phaseportrait, nullclines 
%   and an Euler-path of an inhibited Chemostat.
%   chemostat_inhibited(alfa1, alfa2, np0, cp0, nc) will run if 
%   alpha1 > 1/xc, thus there is a reproduction.
%   alpha2 > 1/(xc*alpha1-1), thus there is sufficient stock-nutrition. 
%   xp0 > 0 , you can not have a nonpositive population.
%   sp0 > 0 , you can not have a nonpositive concentration.
%   xc  > 0
%
%   The blue arrows represent the vectorfield.
%   The black lines are two of the three nullclines.
%   The black dotted line is the invariant line (no solution crosses it).
%   The red line is an Eulerpath, starting in + and ending in *.
%
%   Try the following:
%   chemostat_inhibited(5, 3, 0.2, 0.3, 6)
%
%   by Per Erik Strandberg, 2003-2004.
%

% Start-condition:
%------------------
if ((alpha1>1) & (alpha2>0) & (sp0>0) & (xp0>0) & xc>0),

    if (alpha2<1/(alpha1-1)),
        disp(' ')
        disp (' (HINT: Only the trivial steady state, alpha2 is too small...)')
    else 
        disp(' ')
        disp (' (HINT: Two steady states, alpha2 is quite large...)')
    end
        
    
% The illegal indata case:
%---------------------------
else
    disp('   chemostat_inhibited.m by Per Erik Strandberg, 2003-2004.')
    disp('   Did not Finish OK. (You used illegal indata.)')
    disp(' ')
    disp('   For syntax help type: help chemostat_inhibited .')    
    disp(' ')
end
\end{verbatim}
\end{scriptsize}



% this is the copyright page, I did not touch it much.

\cleardoublepage

\thispagestyle{plain}

\begin{picture}(20,20)(0,0)
% \put(0,10){\includegraphics[width=3.5cm]{epresslogo}}
\put(280,0){\includegraphics[width=2.5cm]{univlogo.eps}}
\end{picture}
 
\null
\vspace{0cm} 

Copyright
\medskip

The publishers will keep this document online on the Internet - or its possible replacement - for a period of 25 years from the date of publication barring exceptional circumstances.
The online availability of the document implies a permanent permission for anyone to read, to download, to print out single copies for your own use and to use it unchanged for any non-commercial research and educational purpose. Subsequent transfers of copyright cannot revoke this permission. All other uses of the document are conditional on the consent of the copyright owner. The publisher has taken technical and administrative measures to assure authenticity, security and accessibility.
According to intellectual property law the author has the right to be mentioned when his/her work is accessed as described above and to be protected against infringement.
For additional information about the Linköping University Electronic Press and its procedures for publication and for assurance of document integrity, please refer to its WWW home page: http://www.ep.liu.se/
\bigskip

Upphovsr\"att
\medskip

Detta dokument hålls tillgängligt på Internet - eller dess framtida ersättare - under 25 år från publiceringsdatum under förutsättning att inga extraordi\-nära omständigheter uppstår.
Tillgång till dokumentet innebär tillstånd för var och en att läsa, ladda ner, skriva ut enstaka kopior för enskilt bruk och att använda det oförändrat för ickekommersiell forskning och för undervisning. Överföring av upphovsrätten vid en senare tidpunkt kan inte upphäva detta tillstånd. All annan användning av dokumentet kräver upphovsmannens medgivande. För att garantera äktheten, säkerheten och tillgängligheten finns det lösningar av teknisk och administrativ art.
Upphovsmannens ideella rätt innefattar rätt att bli nämnd som upphovsman i den omfattning som god sed kräver vid användning av dokumentet på ovan beskrivna sätt samt skydd mot att dokumentet ändras eller presenteras i sådan form eller i sådant sammanhang som är kränkande för upphovsmannens litterära eller konstnärliga anseende eller egenart.
För ytterligare information om Linköping University Electronic Press se förlagets hemsida http://www.ep.liu.se/

\medskip

\copyright\phantom{.} \putshortdate, \putauthor




%----------------END OF DOCUMENT----------------------------------------------
%----------------END OF DOCUMENT----------------------------------------------

\end{document}
 
